\label{lowcollision}\subsection{Fokker-Planck Equation}

The Fokker-Planck Equation is 
\begin{equation}
v_{s}\frac{\partial f_{0}}{\partial s}=\mathcal{C}\left( f_{0}\right) +%
\mathcal{Q}\left( f_{0}\right) +\mathcal{E}\left( f_{0}\right)
\end{equation}

In the low collision regime, which is characterized by the condition
\begin{equation}
\nu ^{*}\equiv \frac{\tau _{b}}{\tau _{dt}}\ll 1
\end{equation}
where $\tau _{dt}=\epsilon \tau _{c}$ is the collision detrapping time, as
defined in the previous section, it is assumed that electrons circulating or
trapped are able to complete their orbit in a time too short for collisions
to deflect them from their orbit. As a consequence, the dominant term in the
Fokker-Planck equation is simply
\begin{equation}
v_{s}\frac{\partial f_{0}}{\partial s}=0
\end{equation}
so that $f_{0}$ is constant along the field lines.

Then, performing a bounce-averaging, we have 
\begin{eqnarray}
\left\{ v_{s}\frac{\partial f_{0}}{\partial s}\right\}  &=&\frac{1}{\tau _{b}%
}\left[ \frac{1}{2}\sum_{\sigma }\right] _{T}\int_{s_{\min }}^{s_{\max }}%
\frac{ds}{\left| v_{s}\right| }v_{s}\frac{\partial f_{0}}{\partial s} 
\nonumber \\
&=&\frac{1}{\tau _{b}}\left[ \frac{1}{2}\sum_{\sigma }\right] _{T}\sigma
\left[ f_{0}\right] _{s_{\min }}^{s_{\max }}  \label{dsbav}
\end{eqnarray}

For passing particles, the positions $s_{\min }$ and $s_{\max }$ coincide,
so that $\left[ f_{0}\right] _{s_{\min }}^{s_{\max }}=0$ and the term
vanishes. For trapped particles, the term also vanishes because of the sum
over $\sigma =\pm 1$, since, by definition, $v_{\Vert }=0$ at the turning
points  $s_{\min }$ and $s_{\max },$ and consequently  $f_{0}$ is
independent of the sign of $\sigma .$ 

Therefore, the bounce-averaged Fokker-Planck equation becomes 
\begin{equation}
\left\{ \mathcal{C}\left( f_{0}\right) \right\} +\left\{ \mathcal{Q}\left(
f_{0}\right) \right\} +\left\{ \mathcal{E}\left( f_{0}\right) \right\} =0
\end{equation}
with $f_{0}$ constant along the field lines.

\subsection{Drift-Kinetic Equation}

The drift kinetic equation is 
\begin{equation}
v_{s}\frac{\partial f_{1}}{\partial s}+\frac{v_{\Vert }}{\Omega _{e}}I\left(
\psi \right) \frac{\left| \nabla \psi \right| }{R}\frac{\partial }{\partial s%
}\left( \frac{v_{\Vert }}{B}\right) \frac{\partial f_{0}}{\partial \psi }=%
\mathcal{C}\left( f_{1}\right) +\mathcal{Q}\left( f_{1}\right) +\mathcal{E}%
\left( f_{1}\right) 
\end{equation}

In the low collisiona regime $\nu ^{*}\ll 1$, the dominant term is 
\begin{equation}
v_{s}\frac{\partial f_{1}}{\partial s}+\frac{v_{\Vert }}{\Omega _{e}}I\left(
\psi \right) \frac{\left| \nabla \psi \right| }{R}\frac{\partial }{\partial s%
}\left( \frac{v_{\Vert }}{B}\right) \frac{\partial f_{0}}{\partial \psi }=0
\end{equation}
which gives 
\begin{equation}
f_{1}=\widetilde{f}+g
\end{equation}
where 
\begin{equation}
\widetilde{f}=-\int ds\frac{v_{\Vert }}{v_{s}\Omega _{e}}I\left( \psi
\right) B_{P}\frac{\partial }{\partial s}\left( \frac{v_{\Vert }}{B}\right) 
\frac{\partial f_{0}}{\partial \psi }
\end{equation}
and $g$ is a constant function along the field lines. Noting that 
\begin{equation}
\widehat{b}\cdot \widehat{s}=-\frac{B_{P}}{B}
\end{equation}
we get 
\begin{eqnarray}
\widetilde{f} &=&\int ds\frac{1}{\Omega _{e}}I\left( \psi \right) B\frac{%
\partial }{\partial s}\left( \frac{v_{\Vert }}{B}\right) \frac{\partial f_{0}%
}{\partial \psi }  \nonumber \\
&=&\frac{\gamma m_{e}I\left( \psi \right) }{q_{e}}\frac{\partial f_{0}}{%
\partial \psi }\int ds\frac{\partial }{\partial s}\left( \frac{v_{\Vert }}{B}%
\right)   \nonumber \\
&=&\frac{v_{\Vert }}{\Omega _{e}}I\left( \psi \right) \frac{\partial f_{0}}{%
\partial \psi }
\end{eqnarray}
where we used the fact that $\partial f_{0}/\partial s=0$.

Then, performing a bounce-averaging, we find again, using the same argument
as in (\ref{dsbav}), that 
\begin{equation}
\left\{ v_{s}\frac{\partial f_{1}}{\partial s}\right\} =0
\end{equation}

In addition, we have 
\begin{eqnarray}
&&\left\{ \frac{v_{\Vert }}{\Omega _{e}}I\left( \psi \right) \frac{\left|
\nabla \psi \right| }{R}\frac{\partial }{\partial s}\left( \frac{v_{\Vert }}{%
B}\right) \frac{\partial f_{0}}{\partial \psi }\right\}   \nonumber \\
&=&\frac{1}{\tau _{b}}\left[ \frac{1}{2}\sum_{\sigma }\right]
_{T}\int_{s_{\min }}^{s_{\max }}\frac{ds}{\left| v_{s}\right| }\frac{%
v_{\Vert }}{\Omega _{e}}I\left( \psi \right) \frac{\left| \nabla \psi
\right| }{R}\frac{\partial }{\partial s}\left( \frac{v_{\Vert }}{B}\right) 
\frac{\partial f_{0}}{\partial \psi }  \nonumber \\
&=&\frac{1}{\tau _{b}}\frac{\gamma m_{e}I\left( \psi \right) }{q_{e}}\frac{%
\partial f_{0}}{\partial \psi }\left[ \frac{1}{2}\sum_{\sigma }\right]
_{T}\sigma \int_{s_{\min }}^{s_{\max }}ds\frac{\partial }{\partial s}\left( 
\frac{v_{\Vert }}{B}\right)   \nonumber \\
&=&\frac{1}{\tau _{b}}\frac{\gamma m_{e}I\left( \psi \right) }{q_{e}}\frac{%
\partial f_{0}}{\partial \psi }\left[ \frac{1}{2}\sum_{\sigma }\right]
_{T}\sigma \left[ \frac{v_{\Vert }}{B}\right] _{s_{\min }}^{s_{\max }}
\end{eqnarray}

Again, for passing particles, the positions $s_{\min }$ and $s_{\max }$
coincide, so that $\left[ v_{\Vert }/B\right] _{s_{\min }}^{s_{\max }}=0$
and the term vanishes. For trapped particles, the term also vanishes because 
$v_{\Vert }\rightarrow 0$ at the turning points $s_{\min }$ and $s_{\max }.$

Consequently, we find that the bounce-averaged drift kinetic equation
becomes 
\begin{equation}
\left\{ \mathcal{C}\left( f_{1}\right) \right\} +\left\{ \mathcal{Q}\left(
f_{1}\right) \right\} +\left\{ \mathcal{E}\left( f_{1}\right) \right\} =0
\end{equation}
where . 
\begin{equation}
f_{1}=\widetilde{f}+g
\end{equation}
\begin{equation}
\widetilde{f}=\frac{v_{\Vert }}{\Omega _{e}}I\left( \psi \right) \frac{%
\partial f_{0}}{\partial \psi }
\end{equation}
and $g$ is then given by 
\begin{equation}
\left\{ \mathcal{C}\left( g\right) \right\} +\left\{ \mathcal{Q}\left(
g\right) \right\} +\left\{ \mathcal{E}\left( g\right) \right\} =-\left\{ 
\mathcal{C}\left( \widetilde{f}\right) \right\} -\left\{ \mathcal{Q}\left( 
\widetilde{f}\right) \right\} -\left\{ \mathcal{E}\left( \widetilde{f}%
\right) \right\} 
\end{equation}
using the fact that all operators are linear.
